\documentclass[10pt,a4paper]{amsart}
\usepackage[margin=19mm]{geometry}
\usepackage{amsmath,amssymb,mathtools}
\usepackage[scr=rsfs,cal=euler]{mathalfa}
\usepackage{xfrac}
\usepackage[unicode,hidelinks,bookmarks=false]{hyperref}
\newcommand{\ie}{i.e.\ }
\newcommand{\cf}{cf.\ }
\newcommand{\resp}{respectively\ }
\newcommand{\Subsection}{\subsection}
\providecommand{\nobreakdash}{\nobreak\hbox{-}}
\setlength{\emergencystretch}{2em}
\multlinegap0pt
\usepackage{bm}
\usepackage{bbm}
\usepackage{dsfont}
\usepackage{xspace}
\usepackage{cancel}
\usepackage{mathtools}
\usepackage{amsthm}
\makeatletter
\@ifundefined{proposition}{\newtheorem{proposition}{Proposition}}{}
\makeatother
\makeatletter
\newcommand{\BE}{D^m}
\newcommand{\DD}{\mathbb{D}}
\DeclareMathOperator{\Hom}{Hom}
\newcommand{\fr}[1]{\mathfrak{#1}}
\@ifundefined{midsecproof}{\newenvironment{midsecproof}[1]{\vspace{\topsep} \noindent \textit{Proof of #1.}}{\hfill$\square$}}{}
\DeclareMathOperator{\op}{op}
\newcommand{\scrE}{\mathcal{E}}
\newcommand{\scrF}{\mathcal{F}}
\newcommand{\scrG}{\mathcal{G}}
\newcommand{\scrH}{\mathcal{H}}
\newcommand{\wistar}{\stackrel{I}{\star}}
\makeatother
\title[Relative Serre duality for Coxeter groups]{Relative Serre duality for Coxeter groups}
\author{Colton Sandvik}
\date{Source: arXiv:2604.06084v1 (CC BY 4.0)}
\begin{document}
\maketitle

\noindent\emph{Source and licence.} Relative Serre duality for Coxeter groups. Version arXiv:2604.06084v1 (CC BY 4.0), distributed under the Creative Commons Attribution 4.0 International licence, as recorded in the arXiv licence metadata for this exact version and shown on the abstract page https://arxiv.org/abs/2604.06084v1. Full rights notice: licenses/arxiv-2604.06084-CC-BY-4.0.txt.\par\smallskip

\noindent\textit{Editorial scope.} Counted unit: the Proposition whose source label is \texttt{prop:pull\_push\_module\_functors} in the article (source0.tex lines 514--528, source label \texttt{prop:pull\_push\_module\_functors}), delivered together with the article's own complete original proof of it (source0.tex lines 529--553, one \texttt{proof} environment of 25 lines). No mathematics is written, repaired or re-derived: statement and proof are the article's own text line for line. Required context retained uncounted: eq:left\_WI\_action (lines 494--496); eq:rigid\_hom\_module (lines 505--507). Every definition, display and label of those lines is retained unchanged. Omitted from this package and not redistributed: 1--493, 497--504, 508--513, 554--783. Nothing inside a retained excerpt is omitted or altered: every retained body is delivered exactly as the article's lines are written, so no omission receipt of any kind accompanies this package. Citations: the retained excerpts carry no citation command at all, so no bibliography entry of the article is transcribed and no reference is redistributed. Reference adaptation: none was needed and none was made; every reference inside the delivered unit resolves inside it, so no unresolved reference or citation is produced. Printed slips kept literal: no printed word, symbol, hypothesis or number of the counted statement or of its proof was altered. Layout and class: the article's own class is \texttt{amsart} with packages including amsmath, amsfonts, amssymb, enumerate, enumitem, hyperref, graphicx, xy, wrapfig, fancyvrb, euscript, fontenc, listings, tikz; the wrapper is the lane's pinned \texttt{amsart} editorial wrapper at 10pt on A4, which restates the theorem-like environments the retained text needs and adds the article's own macro definitions that the retained text needs (\texttt{\textbackslash BE}, \texttt{\textbackslash DD}, \texttt{\textbackslash Hom}, \texttt{\textbackslash fr}, \texttt{\textbackslash op}, \texttt{\textbackslash scrE}, \texttt{\textbackslash scrF}, \texttt{\textbackslash scrG}, \texttt{\textbackslash scrH}, \texttt{\textbackslash wistar}), with extra wrapper packages (bm, bbm, dsfont, xspace, cancel, mathtools). The delivered numbering is the unit's own. Assets: no retained span contains a figure environment, a graphics-inclusion command, a PDF-page inclusion, a table or a TikZ picture; no image file is copied into this package, the package declares no asset dependency and no image file is redistributed with it. Licence: version arXiv:2604.06084v1 of this article is distributed under the Creative Commons Attribution 4.0 International licence, as recorded in the arXiv licence metadata for this exact version and shown on the abstract page https://arxiv.org/abs/2604.06084v1. A full rights notice is delivered at licenses/arxiv-2604.06084-CC-BY-4.0.txt. Characters: every retained excerpt is pure ASCII, so the delivered engine, XeLaTeX with its default Latin Modern text font, reports no missing character.
        \begin{equation}\label{eq:left_WI_action}
                (-) \wistar (-) : \BE (\fr{h}, W_I) \times \BE_X (\fr{h}, W) \to \BE_X (\fr{h}, W) \qquad (\scrF, \scrG) \mapsto (i_X)^* (\iota (\scrF) \star (i_X)_* (\scrG)).
        \end{equation}

        \begin{equation}\label{eq:rigid_hom_module}
                \Hom (\scrF \star^I \scrG, \scrH) \cong \Hom (\scrG, \DD (\scrF)^{\op} \star^I \scrH)
        \end{equation}

        \begin{proposition}\label{prop:pull_push_module_functors}
                Assume that $W$ is finite, and let $X \subseteq Y \subseteq W$ be locally closed subsets. Let $\scrF \in \BE (\fr{h}, W_I)$, $\scrG \in \BE_X (\fr{h}, W)$,  and $\scrH \in \BE_Y (\fr{h}, W)$.
                \begin{enumerate}
                        \item Assume that $X$ and $Y$ are left $W_I$-stable. Then there are natural isomorphisms
                        \begin{align*}
                                \scrF \wistar (i_X^Y)_* \scrG &\cong (i_X^Y)_* (\scrF \wistar \scrG), & \scrF \wistar (i_X^Y)_! \scrG &\cong (i_X^Y)_! (\scrF \wistar \scrG), \\
                                \scrF \wistar (i_X^Y)^* \scrH &\cong (i_X^Y)^* (\scrF \wistar \scrH), & \scrF \wistar (i_X^Y)^! \scrH &\cong (i_X^Y)^! (\scrF \wistar \scrH).
                        \end{align*}
                        \item Assume that $X$ and $Y$ are right $W_I$-stable. Then there are natural isomorphisms
                        \begin{align*}
                                (i_X^Y)_* \scrG \wistar \scrF &\cong (i_X^Y)_* (\scrG \wistar \scrF), & (i_X^Y)_! \scrG \wistar \scrF &\cong (i_X^Y)_! (\scrF \wistar \scrG), \\
                                (i_X^Y)^* \scrH \wistar \scrF &\cong (i_X^Y)^* (\scrH \wistar \scrF), & (i_X^Y)^! \scrH \wistar \scrF &\cong (i_X^Y)^! (\scrH \wistar \scrF).
                        \end{align*}
                \end{enumerate} 
        \end{proposition}

        \begin{proof}
                We will just prove (1) as (2) follows from a symmetric argument.
                First, we will show the isomorphism
                \[ \scrF \wistar (i_X^Y)_* \scrG \cong (i_X^Y)_* (\scrF \wistar \scrG).\]
                The other pushforward isomorphism follows from a similar argument using the earlier observation that $(i_X)_*$ can be replaced by $(i_X)_!$ in the definition of the $\wistar$-bifunctor from (\ref{eq:left_WI_action}).
                We can check the isomorphism from the definitions and using the observation that $\iota (\mathcal{F}) \star (i_X)_* \mathcal{G}$ is in the essential image of $(i_X)_*$.
                \begin{align*}
                        \mathcal{F} \wistar (i_X^Y)_* \scrG &\cong (i_Y)^* (\iota (\scrF) \star (i_X)_* \scrG) \\
                        &\cong (i_X)_* (i_X)^* (i_Y)^*  (\iota (\scrF) \star (i_X)_* \scrG) \\
                        &\cong (i_X^Y)_* (\scrF \wistar \scrG).
                \end{align*}
                
                Next, we will show there is an isomorphism
                \[ \scrF \wistar (i_X^Y)^* (\scrH) \cong (i_X^Y)^* (\scrF \wistar \scrH).\]
                The remaining isomorphism follows from a similar argument.
                Let $\scrE \in \BE_X (\fr{h}, W)$. We can then compute
                \begin{align*}
                        \Hom (\scrF \wistar (i_X^Y)^* \scrH, \scrE) &\cong \Hom ((i_X^Y)^* \scrH, \DD (\scrF)^{\op} \wistar \scrE) \\
                        &\cong \Hom ( \scrH, (i_X^Y)_* ( \DD (\scrF)^{\op}  \wistar \scrE )) \\
                        &\cong \Hom (\scrH,  \DD (\scrF)^{\op} \wistar (i_X^Y)_* \scrE) \\
                        &\cong \Hom (\scrF \wistar \scrH,  (i_X^Y)_* \scrE) \\
                        &\cong \Hom ((i_X^Y)^* (\scrF \wistar \scrH), \scrE).
                \end{align*}
                Here the first and fourth isomorphisms follow from (\ref{eq:rigid_hom_module}) and the third isomorphism follows from the earlier observation that $(i_X^Y)_*$ is left $\BE (\fr{h}, W_I)$-linear.
        \end{proof}

\end{document}
